\documentclass[11pt]{article}
\usepackage[T1]{fontenc}
\usepackage[utf8]{inputenc}
\usepackage[brazilian]{babel}
\usepackage{amsmath,amssymb,amsthm}
\usepackage{newpxtext,newpxmath}
\usepackage[a4paper,margin=24mm,headheight=15pt,headsep=8mm,footskip=11mm]{geometry}
\usepackage{xcolor,booktabs,array,caption,fancyhdr,enumitem}
\usepackage{algorithm,needspace}
\usepackage[noend]{algpseudocode}
\usepackage{tikz}
\usepackage[hidelinks]{hyperref}
\usetikzlibrary{arrows.meta,calc}

\definecolor{ink}{HTML}{223247}
\definecolor{navy}{HTML}{274568}
\definecolor{quiet}{HTML}{617384}
\definecolor{path}{HTML}{C2410C}
\definecolor{ours}{HTML}{2155D6}

\captionsetup{font=small,labelfont={bf,color=navy},labelsep=colon}
\pagestyle{fancy}\fancyhf{}
\renewcommand{\headrulewidth}{.4pt}
\fancyhead[L]{\small\sffamily\color{ink}O CETSP POR GEOMETRIA}
\fancyhead[R]{\small\sffamily\color{quiet}TouringPolygons}
\fancyfoot[L]{\small\color{quiet}28 de outubro de 2026}
\fancyfoot[R]{\small\color{quiet}\thepage}
\setlength{\parindent}{0pt}\setlength{\parskip}{5pt}
\setlist{itemsep=2pt,topsep=3pt}

\theoremstyle{plain}
\newtheorem{lema}{Lema}
\newtheorem{proposicao}{Proposição}
\theoremstyle{remark}
\newtheorem*{obs}{Observação}

\floatname{algorithm}{Algoritmo}
\algrenewcommand\algorithmicrequire{\textbf{Entrada:}}
\algrenewcommand\algorithmicensure{\textbf{Saída:}}
\algrenewcommand\algorithmicfor{\textbf{para}}
\algrenewcommand\algorithmicforall{\textbf{para cada}}
\algrenewcommand\algorithmicdo{\textbf{faça}}
\algrenewcommand\algorithmicwhile{\textbf{enquanto}}
\algrenewcommand\algorithmicif{\textbf{se}}
\algrenewcommand\algorithmicthen{\textbf{então}}
\algrenewcommand\algorithmicelse{\textbf{senão}}
\algrenewcommand\algorithmicreturn{\textbf{devolva}}
\algrenewcommand\algorithmicprocedure{\textbf{procedimento}}
\algrenewcommand\algorithmicfunction{\textbf{função}}
\algrenewcommand{\algorithmiccomment}[1]{\hfill{\color{quiet}\small$\triangleright$ #1}}

\newcommand{\sect}[1]{\par\bigskip{\color{navy}\large\bfseries #1\par}\smallskip}
\newcommand{\R}{\mathbb{R}}
\newcommand{\vez}{\ensuremath{\times}}

\begin{document}

{\color{navy}\LARGE\bfseries O CETSP por geometria\par}
\vspace{2pt}
{\color{quiet}\large o oráculo de discos de ordem fixa e o branch-and-bound sobre ele\par}
\vspace{10pt}

\textbf{O problema.} Dados discos fechados $D_i=\{x\in\R^2:\lVert x-c_i\rVert\le\rho_i\}$,
$i=1,\dots,n$, um ponto inicial $s$ e um ponto final $t$, achar o menor caminho de $s$ a
$t$ que encontra todos os discos. Com $s=t$ (o depósito) é o \emph{close-enough TSP}
clássico, um tour fechado, e é o caso de toda a literatura. Com $s\neq t$ é o análogo,
para discos e com a ordem livre, do TPP de Dror, que liga $s$ a $t$. São duas camadas:
\begin{itemize}
  \item \textbf{Ordem fixa.} Dada uma sequência cíclica $D_{\sigma_1},\dots,D_{\sigma_k}$,
    achar $q_j\in D_{\sigma_j}$ que minimizem $L(q)=\sum_{j=1}^{k}\lVert q_{j+1}-q_j\rVert$, com
    $q_{k+1}=q_1$. O ponto inicial e o final entram como discos de raio $0$. É um problema
    convexo: uma soma de normas com restrições de bola. A literatura o resolve como SOCP;
    aqui ele é resolvido por geometria (seções 1 a 4).
  \item \textbf{Ordem livre.} Um branch-and-bound sobre sequências, com a mesma
    ramificação de Coutinho et al.~(2016). Cada nó chama a ordem fixa (seção 5).
\end{itemize}
Escrevemos $d_j=q_{j+1}-q_j$ e, quando $d_j\neq0$, $u_j=d_j/\lVert d_j\rVert$. Para não
carregar índices, chamamos $D_{\sigma_j}$ de $D_j$, com centro $c_j$ e raio $\rho_j$.

\textbf{Caminho de $s$ a $t$ como ciclo.} O ciclo
$s\to D_{\sigma_2}\to\dots\to D_{\sigma_{k-1}}\to t\to s$ difere do caminho de $s$ a $t$ só
pelo segmento $t\to s$, de comprimento constante $\lVert t-s\rVert$. Os dois têm as mesmas
soluções ótimas, e os limites do ciclo menos $\lVert t-s\rVert$ são limites do caminho. Por
isso o oráculo de ciclos resolve o caminho sem mudança. Com $s=t$, o segmento tem
comprimento zero.

\textbf{A ideia em uma frase.} Como no oráculo de polígonos, o de discos \emph{propõe} um
ciclo por construções geométricas e \emph{prova} limites $L\le\mathrm{OPT}\le U$ com um
certificado que não depende de como o ciclo foi obtido. A proposta é iterativa; a prova é
uma conta em forma fechada.

\sect{1. O certificado}

\begin{proposicao}[limite inferior]\label{prop:dual}
Para quaisquer vetores $u_1,\dots,u_k$ com $\lVert u_j\rVert\le1$, ponha
$w_j=u_{j-1}-u_j$ (com $u_0=u_k$). Então, para todo ciclo $q$ com $q_j\in D_j$,
\[
  L(q)\;\ge\;D(u)\;=\;\sum_{j=1}^{k}\min_{v\in D_j} v\cdot w_j
       \;=\;\sum_{j=1}^{k}\bigl(c_j\cdot w_j-\rho_j\lVert w_j\rVert\bigr).
\]
\end{proposicao}
\begin{proof}
Por Cauchy--Schwarz, $\lVert d_j\rVert\ge u_j\cdot d_j$. Somando e reagrupando por
partes, $\sum_j u_j\cdot(q_{j+1}-q_j)=\sum_j q_j\cdot(u_{j-1}-u_j)=\sum_j q_j\cdot w_j$.
Cada parcela é pelo menos o mínimo de $v\cdot w_j$ sobre o disco, que é atingido em
$v=c_j-\rho_j\,w_j/\lVert w_j\rVert$ e vale $c_j\cdot w_j-\rho_j\lVert w_j\rVert$.
\end{proof}

\textbf{Limite superior.} Qualquer ciclo com $q_j\in D_j$ comprovado dá
$U=L(q)$. O ponto proposto é testado com $\rho_j^2-\lVert q_j-c_j\rVert^2\ge0$ em
arredondamento dirigido; se o teste falha, o ponto é puxado para o centro até passar.

\textbf{Arredondamento.} $D(u)$ e $U$ são avaliados em aritmética intervalar, e
$\lVert u_j\rVert\le1$ é comprovado da mesma forma. Usam-se o extremo inferior de $D$ e o
superior de $U$. Como $\sum_j w_j=0$, os centros podem ser transladados pelo baricentro sem
mudar $D$; isso reduz o cancelamento.

\textbf{Quando o certificado fecha.} Há igualdade na Proposição~\ref{prop:dual} quando:
\begin{itemize}
  \item cada $u_j$ é a direção de $d_j$;
  \item cada $q_j$ minimiza $v\cdot w_j$ no seu disco.
\end{itemize}
Isso tem dois casos:
\begin{itemize}
  \item \emph{passagem reta}: $w_j=0$, o ciclo atravessa o disco sem dobrar;
  \item \emph{reflexão}: $q_j$ está no círculo e $u_j-u_{j-1}$ aponta na normal externa
    em $q_j$. É a lei de reflexão, com ângulo de incidência igual ao de saída.
\end{itemize}
Logo, se a proposta acerta qual disco é reto e qual reflete, e as direções, o certificado
fecha. A correção do resultado depende só destas contas, não da proposta.

\sect{2. Um disco entre dois vizinhos: o problema de Alhazen}

\begin{lema}[atualização de Alhazen]\label{lem:alhazen}
Sejam $a,b\in\R^2$ e $D$ o disco de centro $c$ e raio $\rho>0$. O ponto $x\in D$ que
minimiza $\lVert x-a\rVert+\lVert x-b\rVert$ é:
\begin{itemize}
  \item se o segmento $ab$ encontra $D$: o ponto de $ab$ mais próximo de $c$ (passagem reta;
    qualquer ponto de $ab\cap D$ empata);
  \item senão: o ponto $p(\theta)=c+\rho(\cos\theta,\sin\theta)$ do arco entre as direções
    radiais de $a$ e de $b$ onde $g'(\theta)=0$. Aqui $g(\theta)=\lVert p-a\rVert+\lVert p-b\rVert$,
    $g'(\theta)=\rho\,\tau\cdot(s_a+s_b)$, com $s_a=(p-a)/\lVert p-a\rVert$,
    $s_b=(p-b)/\lVert p-b\rVert$ e $\tau$ a tangente unitária.
\end{itemize}
\end{lema}
Nas duas pontas do arco, a distância ao vizinho daquela ponta é estacionária ao longo do
círculo e a outra diminui para dentro do arco. Então $g'$ troca de sinal no arco, e um
Newton protegido nesse intervalo (com bissecção quando o passo sai) acha a raiz. A raiz é
a lei de reflexão ($s_a+s_b$ paralelo à normal). Em geral é raiz de uma quártica, a
\emph{reflexão de Alhazen}, e por isso não tem forma fechada simples.

\begin{figure}[h]
\centering
\begin{tikzpicture}[scale=0.9,>=Latex]
  \begin{scope}
    \draw[quiet] (0,0) circle (1);
    \fill (0,0) circle (1.2pt) node[below left,quiet,font=\small] {$c$};
    \draw[path,thick] (-2.6,0.4) -- (2.6,-0.2);
    \fill[path] (-2.6,0.4) circle (1.5pt) node[above,font=\small] {$a$};
    \fill[path] (2.6,-0.2) circle (1.5pt) node[above,font=\small] {$b$};
    \fill[ours] (0.02,0.10) circle (1.8pt);
    \node[font=\small,quiet] at (0,-1.5) {passagem reta};
  \end{scope}
  \begin{scope}[xshift=7cm]
    \draw[quiet] (0,0) circle (1);
    \fill (0,0) circle (1.2pt) node[below,quiet,font=\small] {$c$};
    \coordinate (p) at (0,1);
    \draw[path,thick] (-2.2,2.4) -- (p) -- (2.2,2.4);
    \fill[path] (-2.2,2.4) circle (1.5pt) node[left,font=\small] {$a$};
    \fill[path] (2.2,2.4) circle (1.5pt) node[right,font=\small] {$b$};
    \draw[->,ours] (p) -- (0,1.9) node[right,font=\small] {normal};
    \fill[ours] (p) circle (1.8pt) node[below right,font=\small] {$p(\theta)$};
    \node[font=\small,quiet] at (0,-1.5) {reflexão};
  \end{scope}
\end{tikzpicture}
\caption{Os dois casos do Lema~\ref{lem:alhazen}. Na reflexão, os segmentos fazem ângulos
iguais com a normal.}
\end{figure}

\textbf{Relação com os mapas de último passo.} Para um polígono, o mapa de último passo
particiona o plano por vértices e arestas e dá o último ponto de visita por uma fórmula em
cada região. Para um disco, o análogo é a ``sombra'' (onde a passagem é reta) mais um leque
de raios refletidos, parametrizado pelo ângulo $\theta$, monotonamente. A analogia é só
essa: não há partição finita, porque refletir num círculo não é uma isometria, ao contrário
de desdobrar por uma aresta.

\sect{3. O oráculo: proposta geométrica e certificação}

\begin{algorithm}[h]
\caption{Ciclo de ordem fixa por discos (\texttt{tpp\_convex\_solve\_cycle\_disks})}\label{alg:oracle}
\begin{algorithmic}[1]
\Require discos $D_1,\dots,D_k$ na ordem; gap $\varepsilon$; corte $\kappa$ (do B\&B); contatos iniciais opcionais
\Ensure $L\le\mathrm{OPT}\le U$ e contatos $q$ com $L(q)=U$
\State $q_j\gets$ contato inicial ou $c_j$; \Call{Varredura}{$q$} duas vezes
\For{rodada $=1,\dots,32$, enquanto o gap cai à metade em alguma das 3 últimas}
  \State \Call{FecharTipo}{$q$} \Comment{Newton nos ângulos das reflexões}
  \State \Call{Varredura}{$q$} \Comment{confere o tipo}
  \State $u\gets$ \Call{DuaisDoTipo}{$q$};\quad $(L,U)\gets$ certificado (seção 1)
  \If{$U-L\le\varepsilon$ \textbf{ou} $L\ge\kappa$} \Return $(L,U,q)$ \EndIf
\EndFor
\State pontos interiores a partir de $q$: barreira $-\mu\sum_j\log(\rho_j^2-\lVert y_j-c_j\rVert^2)$,
  certificando a cada nível de $\mu$
\State \Return o melhor $(L,U,q)$ certificado
\end{algorithmic}
\end{algorithm}

\textbf{\textsc{Varredura} (Gauss--Seidel).} Cada contato é trocado pela sua atualização de
Alhazen (Lema~\ref{lem:alhazen}) com os vizinhos atuais fixos. Nenhuma troca alonga o
ciclo. É uma busca por blocos de coordenadas com solução exata de cada bloco. Ao fim,
cada disco tem um \emph{tipo}: reto, reflexão, ou \emph{canto} (o ponto onde dois círculos
se cruzam, visitado por dois discos de uma vez).

\textbf{Blocos.} Com discos sobrepostos, dois contatos podem coincidir. O segmento entre
eles tem comprimento zero, e mover um sozinho nunca encurta o ciclo, então a varredura
empaca numa lente. Uma sequência de contatos coincidentes se move junta. Entre os vizinhos
$a$ e $b$, o bloco toma o menor dentre:
\begin{itemize}
  \item um ponto da interseção dos seus discos: no segmento $ab$, se ele a cruza; num ponto
    de Alhazen de um dos círculos, se ele está nos outros; ou num canto;
  \item um membro refletindo e os outros retos;
  \item uma passada de Alhazen em cada sentido, ignorando a coincidência.
\end{itemize}
Um canto só fica se o seu dual o aprova (ver \textsc{DuaisDoTipo}); senão os dois membros
começam um pouco separados e o Newton os separa. É o que acontece nas ``bolhas'' de
Mennell, de círculos tangentes.

\textbf{\textsc{FecharTipo} (Newton).} Com o tipo fixo:
\begin{itemize}
  \item os contatos retos saem do problema, porque ficam nos segmentos entre as dobras;
  \item pontos e cantos ficam fixos;
  \item cada reflexão anda no seu círculo pelo seu ângulo $\theta$.
\end{itemize}
O comprimento do ciclo reduzido é minimizado por Newton nos ângulos. A hessiana é
tridiagonal cíclica e se resolve em $O(k)$ por Thomas com Sherman--Morrison. A curvatura de
cada círculo soma $\rho\lambda\ge0$ à diagonal, com $\lambda$ o multiplicador da reflexão,
então a convergência é quadrática perto do ótimo. Há um deslocamento crescente até o passo
descer, e busca linear.

Se duas dobras se encontram, uma reflexão chegou ao cruzamento do seu círculo com o
seguinte (vira canto) ou a uma dobra fixa (fica ali), e o novo tipo é fechado. Duas
reflexões cujos círculos se cruzam também se fundem no cruzamento quando isso encurta o
ciclo, porque perto de um canto o comprimento não é suave e o Newton fica lento. No fim, os
contatos retos voltam aos seus segmentos, em ordem ao longo de cada um.

\textbf{\textsc{DuaisDoTipo}.} Cada segmento do ciclo recebe a direção do segmento reduzido
que o contém. Assim as passagens retas dobram exatamente zero, e num tipo ótimo fechado
$D(u)=L(q)$ até o arredondamento. Num canto, o segmento de comprimento zero recebe
\[
  u_{\mathrm{meio}}=u_{\mathrm{entra}}+\lambda_1n_1,\qquad
  \lambda_1n_1+\lambda_2n_2=u_{\mathrm{sai}}-u_{\mathrm{entra}},
\]
com $n_1,n_2$ as normais externas dos dois círculos no canto. O canto é ótimo exatamente
quando $\lambda_1,\lambda_2\ge0$ e $\lVert u_{\mathrm{meio}}\rVert\le1$.

\begin{obs}[não há caminho exato]
Para polígonos, um tipo combinatório fixo se resolve exatamente em racionais por
desdobramento, e isso dá o gap zero do Teorema~B do texto sobre o ciclo convexo. Para
discos isso não existe: já um único disco entre dois pontos é uma quártica. O oráculo
certifica $L\le\mathrm{OPT}\le U$ até o gap pedido. Na prática, a construção certifica gaps
de cerca de $10^{-12}$ relativo, e a barreira raramente roda.
\end{obs}

\sect{4. Inserção de um disco: limites em forma fechada}

Seja $q$ o ciclo de um nó, com as suas direções $u$. Insira o disco $E$ na posição $p$, entre
$q_{p-1}$ e $q_p$, e tome o ponto $x^\ast\in E$ mais próximo do segmento $q_{p-1}q_p$. A
Proposição~\ref{prop:dual} com as mesmas direções $u$, exceto que o segmento $q_{p-1}q_p$
dá lugar às direções de $q_{p-1}\to x^\ast$ e de $x^\ast\to q_p$, é um limite inferior
\emph{válido} do filho. Só três parcelas mudam (os dois vizinhos e $E$), então os $k$
limites saem juntos em $O(k)$, sem resolver nenhum filho.

\sect{5. O branch-and-bound}

\begin{algorithm}[h]
\caption{CETSP (\texttt{tpp-convex-cetsp})}\label{alg:bnb}
\begin{algorithmic}[1]
\Require pontos $s$ e $t$ (o depósito, se $s=t$), discos $D_1,\dots,D_n$, gap relativo $\varepsilon=10^{-6}$, tempo
\Ensure caminho factível de comprimento $U^\ast$ e limite inferior $L^\ast$ (comprimentos de ciclo menos $\lVert t-s\rVert$)
\State $U^\ast\gets$ caminho pelos centros; B\&B curto (até $1$\,s) só com ele
\State heurística (até $5\%$ do tempo): inserção mais barata, 2-opt e Or-opt com
  atualização de Alhazen, ILS (inclusive ruína e reconstrução)
\State $Q\gets\{(\langle s\rangle,\,0)\}$ se $s=t$, senão $\{(\langle s,t\rangle,\,0)\}$ \Comment{chave: limite}
\While{$Q\neq\emptyset$}
  \State retire de $Q$ o nó $(\sigma,\beta)$ de menor $\beta$
  \If{$\beta\ge(1-\varepsilon)U^\ast$} \textbf{pare} \EndIf
  \State $(L,U,q)\gets$ Algoritmo~\ref{alg:oracle} em $\sigma$, com corte $\kappa=(1-\varepsilon)U^\ast$
  \If{$\max(\beta,L)\ge\kappa$} \textbf{continue} \Comment{podado} \EndIf
  \State $F\gets\{i:\ D_i\text{ não encontra a poligonal }q\}$ \Comment{sem o segmento $t\to s$}
  \If{$F=\emptyset$}
    \State $U^\ast\gets\min(U^\ast,U)$; o piso da subárvore é $L$ \Comment{tour factível}
  \Else
    \State $E\gets$ o disco de $F$ mais distante de $q$
    \ForAll{posição $p$ de $\sigma$ (se $s\neq t$, só antes de $t$)}
      \State $\beta_p\gets\max(L,\text{limite de inserção de }E\text{ em }p)$ \Comment{seção 4}
      \If{$\beta_p<(1-\varepsilon)U^\ast$} insira $(\sigma+E@p,\,\beta_p)$ em $Q$ \EndIf
    \EndFor
  \EndIf
\EndWhile
\State \Return $U^\ast$ e $L^\ast=\min(\text{limites abertos},\ \text{pisos},\ (1-\varepsilon)U^\ast)$
\end{algorithmic}
\end{algorithm}

Com $s\neq t$, $t$ fica sempre por último, e o gap é medido no comprimento do caminho
(o $\varepsilon$ do algoritmo se aplica a $U^\ast-\lVert t-s\rVert$). A simetria espelhada só
existe com $s=t$.

\textbf{Por que o resultado está certo.} O limite de um nó é o $L$ certificado do ciclo pela
sua sequência. Visitar só parte dos discos nunca alonga o ciclo, então esse limite vale para
toda a subárvore. O $U$ de um tour vem do mesmo certificado.

Detalhes que não mudam a correção:
\begin{itemize}
  \item filhos ``preguiçosos'': entram na fila com o limite de inserção e só são resolvidos
    quando saem;
  \item com $s=t$, a simetria espelhada é removida no segundo disco;
  \item os filhos são guardados em famílias ordenadas;
  \item com um teto de memória, descarta-se o terço de pior limite e guarda-se o menor
    limite descartado, que passa a limitar o que se pode provar;
  \item threads expandem lotes de nós.
\end{itemize}
Os tours finais são conferidos à parte, em aritmética racional, contra o texto decimal das
instâncias.

\textbf{Relação com outros métodos.}
\begin{itemize}
  \item \emph{Coutinho et al.~(2016)}: a ramificação é a mesma. Eles resolvem cada nó com um
    SOCP no CPLEX; aqui entra o Algoritmo~\ref{alg:oracle}, mais os limites de inserção e a
    heurística.
  \item \emph{Busca por blocos de coordenadas} (atualizar um ponto de visita por vez, com
    solução exata de um círculo): é a \textsc{Varredura} do Algoritmo~\ref{alg:oracle}. O
    que se acrescenta é o Newton que fecha o tipo, os blocos para discos sobrepostos e,
    principalmente, o certificado da seção 1. O certificado prova a otimalidade até o gap
    sem hipóteses sobre a convergência da proposta.
\end{itemize}

\needspace{14\baselineskip}
\sect{6. Resultados}

\begin{center}
\small
\renewcommand{\arraystretch}{1.25}
\begin{tabular}{@{}p{5.6cm}p{10.0cm}@{}}
\toprule
Experimento & Resultado \\
\midrule
Ordem fixa \vez\ SOCP do Gurobi, 20 instâncias de Mennell (2D) &
  certifica todas; mediana $17\vez$ mais rápido; o Gurobi acusa erro numérico em 8 \\
B\&B \vez\ Coutinho et al., Tabela 2 (62 instâncias 2D) & fechadas: 32 contra 23 \\
\quad as 23 fechadas por eles &
  mesmos ótimos; $17\vez$ a $18\,000\vez$ mais rápido (mediana $\approx100\vez$), 1 thread \\
\quad as 30 abertas para os dois &
  nosso $L^\ast$ maior em todas: $+2{,}8\%$ a $+20\%$ (mediana $13{,}6\%$); o $L^\ast$ de 4\,h
  deles é superado em $0{,}12$ a $111$\,s (mediana $50$\,s) \\
Behdani e Smith (720 instâncias) & todas fechadas, $3{,}4$\,ms em média \\
$s\neq t$ ($t$ do lado oposto ao depósito) &
  Behdani, 240 com raio $0{,}5$: todas fechadas, $1{,}3\vez$ o tempo de $s=t$ (mediana; pior
  $6\vez$). 22 instâncias rápidas da Tabela 2: mediana $2{,}2\vez$, mas algumas muito mais
  difíceis (rat195-0.10: $0{,}1$ para $43$\,s; d493-0.10 aberta em $50$\,s). Com $t$ perto do
  depósito, só o fator $2$ da simetria perdida. \\
\bottomrule
\end{tabular}
\end{center}

\textbf{Ressalvas.}
\begin{itemize}
  \item As máquinas são de épocas diferentes (um Mac de 2024 contra um i7 de 2012), e a
    diferença não está descontada.
  \item ``Fechada'' quer dizer gap relativo de até $10^{-6}$, com limites provados em
    binary64; não é ótimo algébrico.
  \item Só 2D: a construção é plana. O certificado se estende a esferas; a Tabela 3 (3D) não
    foi tentada.
  \item Os limites inferiores de Carrabs et al.~(2020) ainda não foram comparados.
\end{itemize}

\textbf{No repositório.}
\begin{itemize}
  \item O oráculo: \texttt{packages/convex-tpp/cpp/src/solvers/cycle\_disk\_oracle.cpp}.
  \item O B\&B: \texttt{packages/convex-tpp/cpp/src/main-cetsp.cpp}.
  \item A descrição detalhada: \texttt{docs/algorithms/convex-cycle.md}, seção \emph{Disks}.
  \item Os números: \texttt{benchmarks/results-saved/README.md}, seções \texttt{cetsp-*}.
\end{itemize}

\end{document}
